\documentclass[11pt,a4paper]{article}
\usepackage{fontspec,amsmath,amssymb,amsthm,microtype}
\setmainfont{lmroman10-regular.otf}[BoldFont=lmroman10-bold.otf,ItalicFont=lmroman10-italic.otf,BoldItalicFont=lmroman10-bolditalic.otf]
\usepackage[margin=25mm]{geometry}
\usepackage[unicode,hidelinks]{hyperref}
\setlength{\emergencystretch}{3em}
\setlength{\parskip}{0.4em}
\setlength{\parindent}{0pt}
\allowdisplaybreaks[2]
\raggedbottom
\title{\texorpdfstring{Theorem 1 for Brownian kinetic diffusions}{Theorem 1 for Brownian kinetic diffusions}}
\author{Omer Deniz Akyildiz\\[0.4em]\normalsize Prepared with AI assistance}
\date{}
\begin{document}\maketitle
Draft mathematical exposition; these notes may contain mistakes.


The continuous mixing theorem is instantiated with the existing vector Brownian source. Its initial velocity distribution is arbitrary, subject only to a finite second moment, and is independent of the initial positions and Brownian noise. The output measures are terminal pushforwards of the selected kinetic integral solution, on any finite coordinate set.



The independent-source law theorem is specialized to canonical Brownian forcing. It gives explicit position mixing for the constructed kinetic diffusion with any admissible square-integrable common velocity law.


\section{\texorpdfstring{Construct arbitrary initial velocity and Brownian forcing}{Construct arbitrary initial velocity and Brownian forcing}}


The source is the independent product of the prescribed velocity distribution and the canonical vector Brownian measure. Its first projection gives the L\ensuremath{{}^{2}} initial velocity. The second projection pulls back the canonical Brownian process, multiplied by the fluctuation-dissipation factor \(\sqrt{2\gamma}\). Continuity and the zero initial value follow from the corresponding proved Brownian facts and linear source pullback.



The Brownian factor supplies the physical forcing independently of the prescribed velocity law. Gaussianity of that velocity is unnecessary for the common-input difference estimate. Its finite second moment is enough to construct the phase path in the relevant Hilbert space.


\section{\texorpdfstring{Define the position laws and prove arbitrary-\allowbreak{}shift mixing}{Define the position laws and prove arbitrary-shift mixing}}


For each initial position distribution, the independent-source construction solves the kinetic equation with the original gradient and Brownian input. Its terminal position pushforward defines the law. Applying the general product-source coupling theorem proves the complete arbitrary-shift exponential coefficient. The common velocity law needs no Gaussian assumption. The initial transport term is the true Wasserstein distance between the initial position distributions, obtained by infimizing over couplings rather than assuming an optimizer or a prescribed coupling cost.



The general source-law theorem now specializes to the canonical Brownian input. The complete shifted prefactor remains visible, including any deterioration near the singular coordinate choice. Choosing the simple zero shift yields the familiar exponential kinetic position-mixing bound.


\end{document}
